%% Single source for the Overleaf (LaTeX) and Word versions. Built by build.py.
%% Citations: \cite{key}. Equations: equation environment + \label. Bibliography is generated.

\title{Quantum-assisted allocation of metabolically constrained synaptic plasticity for continual learning in spiking networks}

\author{First Author$^{1}$, Second Author$^{1,2}$ \& Corresponding Author$^{1,*}$}

\affil{$^{1}$Department, Institution, City, Country. $^{2}$Department, Institution, City, Country. $^{*}$e-mail: corresponding.author@institution.edu}

\begin{abstract}
Brains learn continually on a metabolic budget, with unreliable synapses and only local learning signals, whereas neuromorphic hardware lacks a principled rule for deciding which synapses to change when energy is scarce. Here we cast that decision as one constrained optimization problem. Local three-factor plasticity proposes updates; stochastic vesicle release, a metabolic budget and memory consolidation define a quadratic binary allocation problem over dendritic branches that maps onto an Ising Hamiltonian without auxiliary qubits. In simulated spiking networks learning five sequential tasks, budgeted allocation matches the retained accuracy of unconstrained consolidation at 23\% of the plasticity cost. A noiseless simulated depth-two quantum approximate optimization circuit finds the optimal allocation in 90\% of learning windows. We provide circuits for IBM Nighthawk and Heron processors with error suppression, mitigation and detection, a realizable quantum reservoir with matched classical controls and falsifiable criteria for a quantum advantage, which we do not claim.
\end{abstract}

\section*{Introduction}

Synaptic signalling and synaptic modification are metabolically expensive. Most of the energy used by grey matter is spent on restoring ion gradients after synaptic and action-potential activity\cite{attwell2001,harris2012}, and the formation of long-lasting memory carries a cost large enough that animals switch it off when food is scarce\cite{placais2013,li2020}. At the same time, synapses are unreliable: a presynaptic spike releases a vesicle only with some probability, a property that saves energy\cite{levy2002} and that has been linked to Bayesian-like learning\cite{malkin2024,kappel2015} and to the protection of old memories\cite{schug2021}. Learning signals are local. A synapse has access to its own recent activity, stored in an eligibility trace, and to diffuse neuromodulatory or error signals, but not to a global gradient\cite{fremaux2016,bellec2020}. Finally, new learning must not erase old memories, which requires some form of consolidation\cite{fusi2005,benna2016,kirkpatrick2017,zenke2017}.

Neuromorphic computing inherits all four constraints. Event-driven processors are attractive because they spend energy only when spikes occur\cite{roy2019,davies2018}, and on-chip learning relies on local rules that approximate gradients\cite{bellec2020,neftci2019}. Yet writing a synapse is among the most expensive operations on such hardware, and continual operation on a power budget raises a question that current learning rules do not answer: when many synapses are eligible for change and energy allows only a few to be changed, which ones should be chosen? The four constraints are usually studied one at a time. Their interaction is a combinatorial problem, because the value of updating one group of synapses depends on which other groups are updated at the same time.

Quantum computing offers a candidate tool for such problems. Quadratic binary optimization maps directly onto the Ising Hamiltonians that gate-based quantum processors can prepare and sample\cite{lucas2014,farhi2014}, and present superconducting devices execute circuits with thousands of two-qubit gates when error suppression and mitigation are applied\cite{kim2023,ibm2026}. Whether such devices can outperform classical optimization heuristics remains unsettled\cite{abbas2024,koch2026}, and a quantum processor can only help a learning system if the decision it is asked to make is rare, valuable and tolerant of latency. Proposals that place quantum coherence inside neural tissue face strong decoherence objections\cite{tegmark2000} and remain speculative\cite{fisher2015}. We therefore keep a strict separation between the two uses of quantum mechanics in this work: the synapses and neurons of our model are classical, and the quantum processor is an external co-processor that solves a classical allocation problem posed by the network, in the spirit of quantum neuromorphic computing\cite{markovic2020}.

Here we introduce Q-ESPACE (quantum-assisted energy-aware stochastic plasticity and consolidation engine). Its first contribution is a formulation: we show that metabolic cost, release unreliability, local credit assignment and consolidation combine into a single quadratic knapsack problem over dendritic branches, derived from a second-order expansion of the loss, with one binary variable per branch. The second is a qubit-efficient mapping of this problem to an Ising Hamiltonian that needs no slack qubits, together with circuits for the quantum approximate optimization algorithm (QAOA)\cite{farhi2014,blekos2024} matched to the square-lattice IBM Nighthawk and heavy-hex IBM Heron processors\cite{ibm2026}, including error suppression, error mitigation and an error-detecting encoding. The third is a set of controlled simulations. In recurrent spiking networks that learn five tasks in sequence, budgeted allocation retains as much accuracy as unconstrained elastic weight consolidation while spending less than a quarter of the plasticity cost, and a simulated depth-two QAOA sampler finds the optimal allocation in nine of ten learning windows. The fourth is a physically realizable quantum reservoir for multi-timescale credit signals, which we test against matched classical memory and find not to be superior. Throughout, we report negative results alongside positive ones and state the conditions under which a quantum advantage could be claimed. All quantum results in this draft are classical simulations; the hardware experiments are specified and the corresponding entries are marked.

\section*{Results}

\subsection*{Four biological constraints define one allocation problem}

We consider a recurrent network of leaky integrate-and-fire neurons whose plastic output synapses release transmitter stochastically (Fig.~1 and Methods). A presynaptic spike $s_{i,t}\in\{0,1\}$ at synapse $i$ produces a release event
\begin{equation}
r_{i,t}=s_{i,t}\,v_{i,t}\,\xi_{i,t},\qquad \xi_{i,t}\sim\mathrm{Bernoulli}(p_{i,t}),\qquad p_{i,t}=\sigma\!\left(\theta_{0,i}+\theta_{C}C_{i,t}+\theta_{A}\frac{A}{A+K_{A}}\right),
\label{eq:release}
\end{equation}
where $v_{i,t}$ indicates that a vesicle is available, $C_{i,t}$ is a presynaptic calcium trace, $A$ is an ATP-like resource and $\sigma$ is the logistic function. Release events, not spikes, drive the postsynaptic neuron and the eligibility trace $e_{i,t+1}=\lambda_{e}e_{i,t}+r_{i,t}$. With an error signal $m$ delivered at the end of a learning window, the three-factor rule proposes the candidate change
\begin{equation}
\delta_{i}=-\eta\,\big\langle m\,e_{i}\big\rangle_{\mathrm{window}},
\label{eq:threefactor}
\end{equation}
which for output synapses equals a stochastic gradient step\cite{bellec2020}. The candidate is not applied automatically. Synapses are grouped into $n$ dendritic branches\cite{poirazi2020}, and a binary variable $u_{k}$ decides whether branch $k$ commits its candidate, $w\leftarrow w+\sum_{k}u_{k}\delta^{(k)}$, where $\delta^{(k)}$ is $\delta$ restricted to branch $k$. Committing branch $k$ costs $c_{k}$ units of a plasticity resource, and the units available in a window, $B_{t}=\lfloor A_{t}\rfloor$, are drawn from a pool that is replenished at a fixed rate,
\begin{equation}
A_{t+1}=\min\!\big\{A_{\max},\;A_{t}-c^{\top}u_{t}+\rho\big\}.
\label{eq:atp}
\end{equation}

Expanding the window loss $\ell$ to second order in the committed change, adding the increment of an elastic consolidation penalty\cite{kirkpatrick2017} with Fisher importance $F_{i}$ and anchor $w_{i}^{\star}$, a penalty on changes at unreliable synapses and a price on the resource, and using $u_{k}^{2}=u_{k}$, gives
\begin{equation}
f(u)=\sum_{k}a_{k}u_{k}+\sum_{k<l}J_{kl}u_{k}u_{l},\qquad J_{kl}=\delta^{(k)\top}H\,\delta^{(l)},
\label{eq:f}
\end{equation}
\begin{equation}
a_{k}=g^{\top}\delta^{(k)}+\tfrac{1}{2}\delta^{(k)\top}H\,\delta^{(k)}+\frac{\lambda_{F}}{2}\sum_{i\in k}F_{i}\big[2(w_{i}-w_{i}^{\star})\delta_{i}+\delta_{i}^{2}\big]+\lambda_{U}\sum_{i\in k}\delta_{i}^{2}(1-\bar p_{i})+\lambda_{E}c_{k},
\label{eq:a}
\end{equation}
with $g$ and $H$ the gradient and Gauss--Newton curvature of the window loss and $\bar p_{i}$ the mean release probability. The four terms after the gradient term express, in order, interference between updates, consolidation, release unreliability and metabolic price. The decision in each window is then
\begin{equation}
u_{t}^{\star}=\arg\min_{u\in\{0,1\}^{n}}f(u)\quad\text{subject to}\quad c^{\top}u\le B_{t}.
\label{eq:qkp}
\end{equation}
Equation \eqref{eq:qkp} is a quadratic knapsack problem with couplings of either sign, which is NP-hard in general\cite{pisinger2007}. It reduces to an ordinary knapsack problem when $J=0$. Table~1 summarizes how each biological constraint enters.

\subsection*{Mapping the allocation problem to qubits without slack variables}

The standard route from an inequality-constrained problem to an Ising Hamiltonian adds $\lceil\log_{2}(B+1)\rceil$ binary slack variables\cite{lucas2014}. We instead use unbalanced penalization\cite{montanez2024}, which keeps one qubit per branch:
\begin{equation}
F(u)=f(u)-\lambda_{1}h(u)+\lambda_{2}h(u)^{2},\qquad h(u)=B_{t}-c^{\top}u .
\label{eq:unbalanced}
\end{equation}
The penalty is small when the budget is nearly exhausted and grows quadratically when it is exceeded. It does not guarantee feasibility, so every sampled bit string is checked against the original constraint and infeasible samples are discarded; the empty update $u=0$ is always available as a fallback. Substituting $u_{k}=(1-Z_{k})/2$ gives the cost Hamiltonian
\begin{equation}
H_{C}=\sum_{k}h_{k}Z_{k}+\sum_{k<l}K_{kl}Z_{k}Z_{l},\qquad h_{k}=-\frac{a'_{k}}{2}-\frac{1}{4}\sum_{l\neq k}J'_{kl},\qquad K_{kl}=\frac{J'_{kl}}{4},
\label{eq:ising}
\end{equation}
where primes denote the coefficients of $F$. We divide $H_{C}$ by $(\sum_{k}h_{k}^{2}+\sum_{k<l}K_{kl}^{2})^{1/2}$, its standard deviation over uniformly random bit strings, which is available in closed form and makes one set of circuit angles usable across windows. The depth-$p$ QAOA state\cite{farhi2014} is
\begin{equation}
\lvert\psi_{p}(\boldsymbol\gamma,\boldsymbol\beta)\rangle=\prod_{\ell=1}^{p}e^{-i\beta_{\ell}H_{M}}e^{-i\gamma_{\ell}H_{C}}\lvert+\rangle^{\otimes n},\qquad H_{M}=-\sum_{k}X_{k},
\label{eq:qaoa}
\end{equation}
implemented with the gates $R_{Z}(2\gamma_{\ell}h_{k})$, $R_{ZZ}(2\gamma_{\ell}K_{kl})$ and $R_{X}(-2\beta_{\ell})$ (Fig.~2a). The circuit is sampled $S$ times and the best feasible sample is committed. The quantum processor is therefore used as a sampler that concentrates probability on good allocations, and the classical network never waits for a variational loop: the angles are trained once on windows from a development seed and then reused.

\subsection*{Budgeted allocation preserves memory at a fraction of the plasticity cost}

We trained networks with 96 recurrent neurons and $n=12$ branches on five binary spike-pattern tasks presented in sequence, 60 learning windows per task, and measured accuracy on all tasks at the end (20 network seeds; Fig.~3). Unconstrained three-factor learning, which commits every candidate, learned each task well (mean accuracy $0.944$ immediately after training on a task) but retained little: mean final accuracy was $0.678\pm0.011$ (mean $\pm$ s.e.m.) and forgetting was $0.335\pm0.016$. Adding the consolidation penalty to the unconstrained update, the standard elastic-weight-consolidation baseline with its strength selected on development seeds, raised final accuracy to $0.753\pm0.016$ at the same plasticity cost of $7{,}170$ units.

Solving equation \eqref{eq:qkp} exactly in every window gave a final accuracy of $0.757\pm0.013$ and forgetting of $0.184\pm0.018$ while spending $1{,}684$ units, 23.5\% of the unconstrained cost, with on average 3.3 of 12 branches updated per window (Fig.~3c). The accuracy gain over unconstrained learning was $0.078$ (95\% confidence interval $0.061$ to $0.096$, paired $t$-test $P=2\times10^{-8}$); the difference from the consolidation baseline was not significant ($+0.003$, interval $-0.010$ to $0.017$). Choosing a random feasible set of branches at a similar cost gave $0.721\pm0.014$, significantly worse than optimized allocation ($P=4\times10^{-4}$). The benefit thus comes from which branches are chosen, not only from updating fewer of them.

Removing terms from equation \eqref{eq:a} showed which constraints matter in this setting. Without the consolidation term, accuracy fell to $0.690\pm0.013$ ($P=1\times10^{-7}$). Without the unreliability term the change was small ($-0.004$, $P=0.08$). Setting all couplings $J_{kl}$ to zero, which turns the problem into an ordinary knapsack problem, raised the surrogate objective by 9.0\% relative to the optimum but left final accuracy unchanged ($0.756\pm0.014$, $P=0.88$). In this model and at this size, then, interactions between branches affect the quality of the allocation but not the downstream accuracy, which is a negative result for the hypothesis that pairwise terms are needed. We return to its implications below.

\subsection*{Simulated QAOA concentrates probability on optimal allocations}

A noiseless depth-two QAOA sampler with 256 shots per window, driven by one fixed set of four angles, found the optimal feasible allocation in $90.3\pm1.1$\% of windows, with a mean regret of 0.7\% of the optimal objective (Fig.~3d). Under the same protocol simulated annealing reached the optimum in 84.0\% of windows, a greedy rule in 78.2\% and the diagonal knapsack solution in 64.7\%. Drawing 256 uniformly random bit strings and applying the same feasibility filter succeeded in 9.2\% of windows, so the result is not an artefact of post-selection on a small search space. Final accuracy with the QAOA allocator was $0.757\pm0.014$, indistinguishable from exact allocation (difference $0.0001$, interval $-0.004$ to $0.004$).

To separate the sampler from the learning loop, we extracted 40 instances at each of three sizes from networks with different seeds and computed exact output distributions (Fig.~4). With angles trained only on 12-branch development instances and transferred unchanged, the probability that a single shot returns the optimal feasible allocation at depth four was $0.090$ for $n=12$, $0.052$ for $n=16$ and $0.027$ for $n=20$, against $2.4\times10^{-4}$, $1.5\times10^{-5}$ and $9.5\times10^{-7}$ for uniform sampling. The median number of shots needed for a 99\% chance of observing the optimum fell from $1.9\times10^{4}$ (uniform) to 96 at $n=12$ from $3.0\times10^{5}$ to 172 at $n=16$ and from $4.8\times10^{6}$ to 473 at $n=20$. The slack-variable encoding, which needs 15 qubits on average for the 12-branch problem, gave a single-shot success probability of $0.017$ at depth two, less than half of the $0.039$ obtained with 12 qubits by unbalanced penalization.

These instances are small. At $n=20$ the best of 256 shots at depth four had a regret of 1.4\%, lower than the greedy rule (2.0\%) and than simulated annealing with a 2,000-step schedule (3.8\%). This is not an advantage: exhaustive enumeration takes about a millisecond at $n=12$ and 0.64 s at $n=20$ on one processor core, and five annealing runs of 20,000 steps solved every 20-branch instance exactly in 0.29 s. The simulations therefore establish that the mapping and the fixed-angle protocol work, and say nothing about a speed-up.

\subsection*{Noise, mitigation and error detection}

We next simulated the circuits under a depolarizing gate-noise and read-out model with error rates of the order reported for current Heron-class devices\cite{ibm2026} (Methods). Dense 10-branch instances were charged three two-qubit gates per interaction to account for routing, giving 270 two-qubit gates at depth two. The two outputs of the circuit respond differently to noise (Fig.~5). The expectation value $\langle H_{C}\rangle$, which is the quantity used to train angles, was biased by $13.0$\% at a two-qubit error of $2.5\times10^{-3}$. Read-out mitigation\cite{vandenberg2022} reduced this to $11.8$\% and zero-noise extrapolation\cite{temme2017,li2017} with an exponential fit to $2.8$\%. The samples, which are what the network actually uses, degraded gracefully: the single-shot probability of the optimal allocation fell from $0.067$ to $0.044$, and the regret of the best of 256 shots was $0.5$\%, because one good sample among many suffices and feasibility is verified classically.

Fault-tolerant error correction is beyond the resources of present devices for this workload, but error detection is available. In the $[[k+2,k,2]]$ code\cite{self2024} every QAOA gate becomes a two-body rotation (Fig.~2c), and a shot is discarded when the measured parity is odd. In simulation with the same error rate per gate and one two-qubit gate per interaction, the encoded circuit kept $77$\% of shots at a two-qubit error of $2.5\times10^{-3}$ and its accepted shots contained the optimal allocation with probability $0.066$, against $0.061$ for the unencoded circuit; at an error of $10^{-2}$ the figures were $0.039$ and $0.042$ with $57$\% of shots kept. Detection therefore gave a small gain at low error rates ($P=0.05$ at $2.5\times10^{-3}$, 10 instances) and none at $10^{-2}$, where the encoded circuit, with 141 two-qubit gates against 90, accumulates more errors than the parity check removes. This comparison ignores connectivity. On a lattice the two code qubits that touch every data qubit require additional routing, so the benefit must be established on hardware.

\subsection*{Circuits and resources for IBM Nighthawk and Heron processors}

Two properties of the problem determine the hardware cost: the number of branches and the density of couplings. For dense couplings we use an odd--even network of nearest-neighbour swaps on a line of qubits (Fig.~2b), in which every pair of branches interacts once per QAOA layer\cite{weidenfeller2022}; a line of the required length exists on both the heavy-hex lattice of Heron and the square lattice of Nighthawk. Its two-qubit gate count grows as $\tfrac{3}{2}n(n-1)p$ (Table~2). At depth two a 12-branch problem needs 396 two-qubit gates and a 40-branch problem 4,680, close to the 5,000-gate capability stated for the first Nighthawk processor\cite{ibm2026}. For larger networks the couplings must be sparse. If branches are placed on a square grid and only couplings between neighbouring branches are kept, the problem is native to the Nighthawk lattice: a $10\times12$ grid that uses all 120 qubits needs 872 two-qubit gates at depth two, with a two-qubit depth of 16. In our instances this truncation kept 17 to 25\% of the total coupling weight and increased the objective by 5.5\% at $n=12$, 15.3\% at $n=16$ and 16.5\% at $n=20$ relative to the optimum, compared with 10.3\%, 16.5\% and 21.8\% when all couplings were dropped. Truncation thus recovers part of the value of the couplings, and how much is lost in accuracy must be measured for each network.

The accompanying code builds each circuit, transpiles it for a named backend and applies dynamical decoupling\cite{viola1999}, Pauli twirling\cite{wallman2016}, twirled read-out mitigation and zero-noise extrapolation through the Qiskit Runtime primitives\cite{qiskit2024}. Because the first Nighthawk release did not support dynamic circuits\cite{ibm2026}, sampling experiments are assigned to Nighthawk and experiments that need mid-circuit measurement to a Heron r3 processor. \todo{Hardware results. Insert, for each backend: transpiled two-qubit gate count and depth, single-shot probability of the optimal allocation, feasible fraction, regret of the best of 256 shots, mitigated and unmitigated $\langle H_{C}\rangle$ against the noiseless value, accepted fraction for the error-detected circuit, and wall-clock and QPU time per window. The scripts that produce these numbers are described in Methods.}

\subsection*{A physically realizable quantum reservoir does not outperform matched classical memory}

Eligibility traces span several timescales, and open quantum systems with memory have been proposed as temporal feature generators\cite{fujii2017,sannia2026}. A reservoir that mixes its present state with a stored copy of a past state cannot be built, because unknown quantum states cannot be copied. We therefore use a collision model\cite{ciccarello2022} (Fig.~2d): three memory qubits interact at every time step with two interface qubits that are reset, loaded with the inputs, entangled with the memory and measured\cite{mujal2023}. Only the memory qubits carry information forward, and every operation is available on processors with mid-circuit measurement.

Four such five-qubit reservoirs provide 12 features. In exact simulation they had a linear memory capacity of $3.59\pm0.07$ (10 input sequences), which fell to $2.28\pm0.06$ with 2,048 shots per feature (Fig.~6a). An echo-state network\cite{jaeger2004} with 12 units and the same hyperparameter search budget reached $5.89\pm0.23$, and a 12-tap delay line reached the maximum of 9 for the delays tested. On a nonlinear task, the delayed product $s_{t}s_{t-\tau}$, the normalized test error was $0.267\pm0.009$ for the quantum reservoir, $0.124\pm0.008$ for the echo-state network and $0.142\pm0.004$ for the delay line. On the task motivated by this work, reconstructing the three-factor credit signal $m_{t}\sum_{k}\lambda^{k}s_{t-1-k}$ from the two input streams, the errors were $0.110\pm0.006$, $0.057\pm0.009$ and $0.045\pm0.002$, and a read-out given the exact product terms solved the task. The quantum reservoir thus has genuine fading memory that survives 5\% dephasing per step, but for classical input streams at this size it offers no benefit over classical memory with the same number of features.

\section*{Discussion}

We asked how quantum computing can be connected to the constraints under which brains and neuromorphic systems learn without assuming that neural tissue is itself quantum. The answer proposed here is that the constraints jointly define a discrete decision, and that this decision has the mathematical form quantum optimizers are built for. Metabolic cost supplies the budget, consolidation and unreliability shape the linear terms, and interference between simultaneous updates supplies the couplings. Posed this way, the plasticity decision can be handed to any solver, and comparing solvers becomes an empirical matter.

Three findings follow from the simulations. First, deciding where to spend plasticity is worth doing: optimized allocation retained as much as unconstrained consolidation while writing less than a quarter as much, and random allocation at the same cost was worse. For neuromorphic devices in which synaptic writes dominate energy and wear, this is a direct lever, and it does not depend on the quantum component. Second, the allocation problem can be sampled by a shallow quantum circuit with one qubit per branch and fixed angles. The fixed-angle protocol matters in practice, because it removes the variational loop from the learning path and lets the processor be called asynchronously, once per consolidation window, at a rate set by the plasticity budget and not by spike times. Third, the simulations expose the limits of the approach as clearly as its promise. At 12 to 20 branches every classical method we tried is fast and accurate, the couplings that make the problem hard did not change downstream accuracy, and a quantum reservoir with physically legitimate memory was outperformed by a delay line.

These negative results sharpen what would have to be true for a quantum advantage. The instances must be large enough to defeat exhaustive and heuristic classical search, which in our formulation means hundreds of coupled branches; the couplings must matter for learning, which is more likely in deep or recurrent plastic networks where curvature couples distant synapses than in the single plastic layer studied here; and the full cost of the quantum call, including compilation, queueing, repetition and classical post-processing, must be lower than that of the best classical solver at the same solution quality\cite{abbas2024,koch2026}. We propose three tests in order of strength: equal or better regret at a fixed time budget; shorter time to a target objective with 99\% confidence, measured end to end; and lower measured energy for the same retained accuracy. None is met by the present work, and the first can be examined on current hardware with the lattice-native instances described above.

Several limitations should be stated plainly. Plasticity was confined to output synapses of a fixed recurrent network, tasks were synthetic, and energy was counted in dimensionless units that have not been calibrated against ATP consumption or device power. The second-order expansion is local, so its predictions were not checked against the realized loss change beyond the window. Noise was modelled as depolarizing and uncorrelated, which omits crosstalk and leakage. The error-detecting code was compared without routing overhead. All quantum data are simulated. Each of these can be addressed with the released code, and the hardware protocol is specified so that the missing measurements can be added without changing the analysis.

More broadly, the formulation suggests a division of labour between physical substrates. Spiking hardware is well suited to fast, local, event-driven computation, and it is poorly suited to global combinatorial choices. Quantum processors are slow and remote, and they are designed for exactly such choices. The biological analogy is loose but suggestive: fast synaptic events propose changes, and slower, resource-limited processes decide which of them are consolidated\cite{frey1997}. Whether a quantum co-processor is the right device for the slow decision is an open and testable question, and the framework presented here states the conditions under which the answer would be yes.

\section*{Methods}

\subsection*{Spiking network and synapse model}
The network has $N_{\mathrm{in}}=60$ input channels and $N=96$ recurrent leaky integrate-and-fire neurons simulated for $T=40$ time steps per sample. Membrane potentials follow $V_{j,t+1}=\alpha_{m}V_{j,t}+I_{j,t}-z_{j,t}V_{\mathrm{th}}$ with $z_{j,t}=\mathbb{1}[V_{j,t}\ge V_{\mathrm{th}}]$, $\alpha_{m}=0.85$ and $V_{\mathrm{th}}=1$; input and recurrent weights are fixed, random and sparse (connection probabilities 0.25 and 0.10). Each recurrent neuron makes one plastic synapse onto a leaky output unit. Release follows equation \eqref{eq:release} with calcium trace $C_{i,t+1}=\alpha_{C}C_{i,t}+\kappa_{C}s_{i,t}$ and a releasable pool
\begin{equation}
R_{i,t+1}=\mathrm{clip}_{[0,1]}\!\Big[R_{i,t}+\frac{1-R_{i,t}}{\tau_{R}}-\nu_{R}r_{i,t}\Big],\qquad v_{i,t}\sim\mathrm{Bernoulli}(R_{i,t}),
\label{eq:vesicle}
\end{equation}
with $\alpha_{C}=0.8$, $\kappa_{C}=0.5$, $\theta_{C}=0.8$, $\theta_{A}=0.5$, $K_{A}=5$, $\tau_{R}=8$, $\nu_{R}=0.35$ and $\theta_{0,i}\sim\mathcal N(0.4,0.6^{2})$. About half of all spikes produced a release. The ATP term in equation \eqref{eq:release} was held at its resting value for signalling; the pool of equation \eqref{eq:atp} ($A_{\max}=12$, $\rho=6$ per window) limits plasticity only. The eligibility trace ($\lambda_{e}=0.92$) at the end of a sample, centred and scaled by statistics measured once under random input, is the feature $x_{i}$ seen by the output unit, whose prediction is $\pi=\sigma(w^{\top}x+b)$. The 96 synapses are divided into $n=12$ branches of eight, with integer costs $c_{k}$ drawn uniformly from $\{1,2,3\}$. All energies are dimensionless.

\subsection*{Tasks and training protocol}
Each task is a binary classification of spike patterns: each class activates seven randomly chosen input channels at rate 0.30 per step during the first or second half of the sample, on a background rate of 0.04 with per-sample rate jitter. Five tasks are presented in sequence, each for 60 windows of 10 samples; one output unit is shared by all tasks. In each window the gradient $g=\langle(\pi-y)x\rangle$, curvature $H=\langle\pi(1-\pi)xx^{\top}\rangle$ and candidates $\delta=-\eta g$ ($\eta=1$) are computed, the allocation problem is solved, the chosen branches are committed and the pool is updated. After each task the diagonal Fisher estimate $F_{i}=\langle\pi(1-\pi)x_{i}^{2}\rangle$ is accumulated and the anchor $w^{\star}$ is set to the current weights. We used $\lambda_{F}=8$, $\lambda_{U}=0.5$ and $\lambda_{E}=2\times10^{-4}$. Hyperparameters, including the consolidation strength of the unconstrained baseline ($\lambda_{F}=2$) and the penalty weights, were chosen on development seeds (100 to 102 and 991) that were not used for the reported results. Spike trains and release events were generated once per seed and shared by all allocators. Forgetting is the mean over tasks of the highest accuracy reached on a task minus its final accuracy.

\subsection*{Constraint encodings and classical solvers}
For unbalanced penalization (equation \eqref{eq:unbalanced}) we set $\lambda_{1}=0.5\,\bar a$ and $\lambda_{2}=0.01\,\bar a$, with $\bar a$ the mean absolute linear coefficient of the instance, giving $a'_{k}=a_{k}+\lambda_{1}c_{k}+\lambda_{2}(c_{k}^{2}-2Bc_{k})$ and $J'_{kl}=J_{kl}+2\lambda_{2}c_{k}c_{l}$. The slack encoding adds $P(c^{\top}u+\sum_{h}2^{h}b_{h}-B)^{2}$ with $\lceil\log_{2}(B+1)\rceil$ slack bits. Classical comparators were exhaustive enumeration, the dynamic-programming solution of the diagonal problem, a greedy rule that repeatedly adds the branch with the largest marginal gain per unit cost, simulated annealing over single-bit flips restricted to feasible states (1,500 to 2,000 steps, geometric schedule), a random feasible mask, and uniform random bit strings passed through the same post-processing as quantum samples. Regret is $(f(u)-f^{\star})/|f^{\star}|$ with $f^{\star}$ the exact feasible optimum; in the learning experiment it is aggregated as a ratio of sums over windows.

\subsection*{QAOA simulation and angle transfer}
Output distributions were computed by state-vector simulation. Angles were trained by minimizing the mean normalized $\langle H_{C}\rangle$ over 12 development instances with the Nelder--Mead method, increasing the depth one layer at a time and initializing each depth by interpolating the previous solution\cite{zhou2020}. The same angles were used for all test instances and, in the solver benchmark, for all sizes. In each window $S=256$ samples were drawn, infeasible samples were discarded, and the feasible sample with the lowest $f$ was committed if $f<0$. The number of shots for a 99\% chance of observing the optimum is $\ln(0.01)/\ln(1-p_{\mathrm{opt}})$, where $p_{\mathrm{opt}}$ is the single-shot probability of an optimal feasible bit string.

\subsection*{Noise model and error mitigation}
Noisy circuits were simulated by averaging 150 Pauli trajectories. Each two-qubit gate was followed with probability $p_{2}$ by a uniformly random non-identity two-qubit Pauli operator and each physical one-qubit rotation with probability $p_{1}=2\times10^{-4}$ by a random one-qubit Pauli operator; $R_{Z}$ was treated as virtual. Read-out flipped 0 to 1 with probability 0.008 and 1 to 0 with probability 0.012. For dense problems the error probability per interaction was $1-(1-p_{2})^{3}$. Read-out mitigation inverts the single-qubit assignment maps on one- and two-body correlators: with $s_{k}=1-p_{01}-p_{10}$ and $o_{k}=p_{10}-p_{01}$, $\langle Z_{k}\rangle=(\langle Z_{k}\rangle_{\mathrm{meas}}-o_{k})/s_{k}$ and correspondingly for $\langle Z_{k}Z_{l}\rangle$. Zero-noise extrapolation evaluated $\langle H_{C}\rangle$ at noise factors 1, 3 and 5 with 4,096 shots each and extrapolated to zero with linear, quadratic and exponential fits; the exponential form is appropriate because the normalized $H_{C}$ is traceless.

\subsection*{Error detection}
The $[[k+2,k,2]]$ code has stabilizers $X^{\otimes(k+2)}$ and $Z^{\otimes(k+2)}$ and logical operators $\bar X_{i}=X_{t}X_{i}$ and $\bar Z_{i}=Z_{i}Z_{b}$, where $t$ and $b$ are the two additional qubits\cite{self2024}. The logical state $\lvert+\rangle^{\otimes k}$ is a GHZ state in the $X$ basis. QAOA gates become $e^{-i\theta\bar Z_{i}}=e^{-i\theta Z_{i}Z_{b}}$, $e^{-i\theta\bar Z_{i}\bar Z_{j}}=e^{-i\theta Z_{i}Z_{j}}$ and $e^{-i\theta\bar X_{i}}=e^{-i\theta X_{t}X_{i}}$\cite{he2024}. After measuring all qubits in the computational basis, shots with odd parity are rejected and logical bits are $u_{i}=m_{i}\oplus m_{b}$. We verified numerically that the noiseless encoded circuit reproduces the unencoded distribution exactly. The hardware circuit optionally measures both stabilizers mid-circuit with two ancilla qubits.

\subsection*{Hardware protocol}
Circuits are generated with Qiskit\cite{qiskit2024} and executed through the Qiskit Runtime Sampler and Estimator primitives. Dense instances use the swap-network circuit and lattice-native instances the direct circuit; both are transpiled at the highest optimization level. Sampler jobs enable dynamical decoupling (XpXm sequence) and gate and measurement twirling; Estimator jobs additionally enable twirled read-out mitigation and zero-noise extrapolation with noise factors 1, 3 and 5. Sampling experiments target the 120-qubit square-lattice Nighthawk processor (\texttt{ibm\_miami}); the error-detected circuits with syndrome measurement and the reservoir target a 156-qubit Heron r3 processor (\texttt{ibm\_boston}). For each job the script stores raw counts, the transpiled gate counts and depth, the physical qubits used, job identifiers and timing. \todo{Report backend names, calibration dates, median two-qubit and read-out errors on the qubits used, shots per circuit, number of twirling randomizations, and queue, QPU and wall-clock times.}

\subsection*{Quantum reservoir and classical baselines}
Each reservoir has three memory qubits with state $\rho^{M}_{t}$ and two interface qubits. At step $t$ the interface qubits are prepared in $R_{Y}(\pi s_{t})\lvert0\rangle$ and $R_{Y}(\pi m_{t})\lvert0\rangle$ (the second in $\lvert0\rangle$ for single-input tasks), all five qubits evolve under $U=e^{-iH\tau}$ with $H=\sum_{i}a_{i}X_{i}+\sum_{i<j}J_{ij}Z_{i}Z_{j}$, $a_{i}\sim\mathcal U(0.5,1.5)$, $J_{ij}\sim\mathcal U(-1,1)$, and the features are $\langle Z\rangle$ of each interface qubit and their correlator. The memory state is updated by the partial trace
\begin{equation}
\rho^{M}_{t+1}=\mathrm{Tr}_{I}\!\big[U\big(\rho^{M}_{t}\otimes\lvert\phi(s_{t},m_{t})\rangle\langle\phi(s_{t},m_{t})\rvert\big)U^{\dagger}\big].
\label{eq:collision}
\end{equation}
Four reservoirs with different interaction times $\tau$ give 12 features. Inputs were uniform on $[0,1]$, sequences had 1,500 steps, and a ridge read-out was trained on steps 50 to 1,000 and tested on the remainder. Memory capacity is $\sum_{\tau=0}^{8}r^{2}_{\tau}$, with $r^{2}_{\tau}$ the squared correlation between the read-out and $s_{t-\tau}$. The echo-state network had 12 tanh units; six settings of interaction times for the quantum reservoir and six settings of spectral radius and leak for the echo-state network were compared on a validation split. The delay line used 12 taps divided among the inputs, and the oracle added the exact nonlinear terms of each task. Finite sampling was modelled by binomial noise with 2,048 shots per feature and dephasing by a phase flip with probability 0.05 per memory qubit per step.

\subsection*{Statistics and reproducibility}
The learning experiment used 20 network seeds, the solver benchmark 40 test instances per size from two seeds, the noise study 16 instances and the reservoir study 10 input sequences. Values are mean $\pm$ s.e.m. unless stated. Comparisons between allocators are paired by seed and use two-sided paired $t$-tests without correction for multiple comparisons; the primary comparisons (optimized allocation against unconstrained learning, against the consolidation baseline and against random allocation, and QAOA against exact allocation) were specified in the analysis script before the test seeds were run. Ten unit tests check the encodings, the Ising map, the agreement between the gate-level and state-vector simulators, the equivalence of encoded and unencoded circuits, the swap network and the read-out correction.

\section*{Data availability}
The instances, simulation outputs and figure source data generated in this study are provided with the code. \todo{Deposit in a public repository and insert the DOI. Add hardware job data when available.}

\section*{Code availability}
The simulation code, experiment scripts, tests and the Qiskit programs for IBM Quantum processors are available at \todo{repository URL and archived DOI}.

@REFERENCES@

\section*{Acknowledgements}
\todo{Funding sources and grant numbers; access to IBM Quantum services if applicable, with the required IBM acknowledgement text.}

\section*{Author contributions}
\todo{Describe each author's contribution using initials.}

\section*{Competing interests}
The authors declare no competing interests. \todo{Confirm or amend.}

@DISPLAY@
